\section{Capítulo 2: Estado del Arte}
En esta sección se presentan investigaciones y proyectos previos, tanto a nivel
nacional como internacional, así como trabajos desarrollados dentro del Instituto
Politécnico Nacional (IPN), en la carrera de Telemática de UPIITA, que están
directamente relacionados con el prototipo propuesto. Estos estudios y proyectos
destacan la relevancia del uso de sensores para la recolección de datos y la
aplicación de algoritmos de inteligencia artificial en el análisis de estos datos, con el
fin de optimizar la calidad y eficiencia en procesos de fermentación o en el monitoreo
de alimentos. Las soluciones revisadas en este contexto están estrechamente
vinculadas con los módulos propuestos en el sistema de monitoreo, particularmente
en lo que respecta a la adquisición de datos y su posterior análisis.
\subsection{Internacional}
\subsubsection{Sistema Embebido para el seguimiento y control en línea de un
fermentador aplicado a la conversión de azúcares en etanol
}
\label{sec:In1}
En este proyecto, se desarrolló un sistema embebido basado en Arduino y NodeRED, diseñado para supervisar y ajustar en tiempo real parámetros críticos como
temperatura y pH en un biorreactor. Permitiendo el seguimiento remoto del proceso
fermentativo mediante una aplicación móvil y visualización web, mejorando así la
eficiencia, la reproducibilidad y la flexibilidad operativa. Validando dicho proceso en
fermentaciones con kéfir, destacando su utilidad para controlar procesos con
levaduras y bacterias en aplicaciones alimentarias o bioenergéticas. Destacando
que el uso de sistemas embebidos es una herramienta clave para el monitoreo y
control de fermentaciones \cite{Paredes}.

\subsubsection{Modelado de la producción de kéfir con lógica difusa}
\label{sec:In2}
Este estudio propuso un modelo de lógica difusa para optimizar la producción
industrial de kéfir, evaluando variables como temperatura de incubación, tiempo de
fermentación y proporción de inóculo. Empleando reglas Mamdani y funciones de
pertenencia triangulares, se construyó un sistema capaz de predecir si la
fermentación alcanzaría el valor de pH ideal (4.6). La lógica difusa redujo la
dependencia del operador y el tiempo de prueba-error, mejorando la eficiencia y
reduciendo costos. Este modelo demostró que es posible automatizar y controlar
procesos fermentativos complejos sin requerir modelos matemáticos precisos, lo
cual es especialmente útil en sistemas alimentarios con variabilidad biológica \cite{Akgul2014}.

\subsubsection{Diseño de fermentador de yogur con sistema de control lógico difusa
usando el microcontrolador ATMEGA32}
\label{sec:In3}
Este proyecto desarrolló un fermentador de yogurt equipado con un sistema de
control basado en lógica difusa implementado en un microcontrolador ATMega32.
El objetivo fue estabilizar la temperatura de fermentación entre 40°C y 45°C, rango
óptimo para el crecimiento de Lactobacillus y Streptococcus. A diferencia del control
ON-OFF tradicional, el sistema con lógica difusa ofreció un perfil térmico más
estable, reduciendo errores de medición y mejorando la calidad del producto final.
El yogurt resultante mostró valores óptimos de pH (3.66) y una mayor concentración
de bacterias, evidenciando la efectividad del enfoque difuso en sistemas no lineales
y biológicos como la fermentación [16].
\subsection{Nacional}
\subsubsection{Sistema de aprendizaje automático para la predicción del avance de
fermentación en cervecería artesanal}
\label{sec:Na1}
Este proyecto, realizado en colaboración con la cervecería artesanal Loba en
Guadalajara, Jalisco, implementó un sistema basado en la nube para monitorear la
temperatura durante la fermentación de la cerveza. Haciendo uso de una IA, el
sistema predice la densidad del mosto y clasifica las etapas de fermentación,
mejorando la consistencia entre lotes y reduciendo pérdidas por variaciones no
detectadas de temperatura. Demostrando ser una solución accesible y efectiva para
pequeñas cervecerías artesanales [17].

\subsubsection{Diseño de un proceso automatizado de fermentación de tepache}
\label{sec:Na2}
En el Tecnológico Nacional de México, se diseñó un prototipo automatizado para la
fermentación de tepache utilizando kéfir. El sistema controla la fermentación
mediante el monitoreo del pH y la regulación de la temperatura, deteniendo el
proceso automáticamente al alcanzar los parámetros establecidos. Este enfoque
busca estandarizar la producción del tepache, mejorando su calidad y seguridad
alimentaria, y facilitando su producción a mayor escala [18].

\subsubsection{Desarrollo de un sensor de $CO_{2}$ y etanol para el monitoreo en
procesos biológicos empleando una plataforma de código abierto}
\label{sec:Na3}
Investigadores de la Universidad Autónoma Metropolitana, unidad Azcapotzalco,
desarrollaron un prototipo de sensor basado en una plataforma de código abierto
para monitorear en tiempo real las concentraciones de $CO_{2}$ y etanol durante
procesos de fermentación. En el sistema se implementó una tarjeta Arduino con
microcontrolador Atmega 328 y sensores electroquímicos, ofreciendo una
alternativa económica para el monitoreo en bioprocesos. Demostrando su eficacia
en la detección de cambios en la concentración de gases, permitiendo un
seguimiento preciso del proceso fermentativo [19].

\subsubsection{Análisis de la innovación de proceso y producto de un sistema de
control y monitoreo de la fermentación alcohólica en vino}
\label{sec:Na4}
Este estudio, realizado en CETYS Universidad, evaluó la viabilidad de implementar
un sistema de control para monitorear la fermentación del vino. Analizando
parámetros clave en el proceso de calidad y se propusieron mejoras tecnológicas
para optimizar y reducir costos de producción. Concluyendo que la automatización
del monitoreo de la fermentación puede mejorar la calidad del vino y aumentar la
competitividad de las bodegas mexicanas en el mercado nacional e internacional
[20].

\subsection{Institucional IPN-UPIITA}
\subsubsection{Sistema Telemático para el monitoreo y control de huertos urbanos
basado en una red inalámbrica de sensores}
\label{sec:IPN1}
En este trabajo realizado por estudiantes de telemática de la UPIITA, se desarrolló
un sistema telemático de monitoreo para el control de huertos, implementando 3
módulos. El primer módulo encardado de la red de sensores y actuadores. El
segundo modulo encargado del servidor y de la base de datos, y por último el tercer
módulo siendo la página web donde el usuario da seguimiento al estado de los
huertos por medio de una interfaz gráfica. Destacando en este trabajo la importancia
del monitoreo en los huertos para mantener un control de calidad en los cultivos
finales de estos [21]. 
\subsubsection{Prototipo de sistema inálambrico de monitoreo y advertencia para visibilizar fugas de mónoxido de carbono en un entorno limitado mediante técnicas de realidad aumentada}
\label{sec:IPN2}
\subsection{Tabla comparativa}
A continuación, se presenta la \cref{table:proyec} cual muestra tanto las principales
características, así como los módulos en donde presentan similitudes marcados con
un \checkmark, los proyectos, investigaciones y estudios antes mencionados con el sistema.

\vspace{0.5cm}
% Add these lines to your document preamble (before \begin{document}):
% \usepackage{multirow}
% \usepackage{array}

\begingroup
\setlength{\tabcolsep}{3pt} % Espacio entre columnas ajustado
\renewcommand{\arraystretch}{1.3} % Altura de fila cómoda
\normalsize % Usamos tamaño normal (o \small si el tamaño normal es muy grande)
% Anchos de columna ajustados para caber en el ancho vertical (aprox. 16cm total)
\begin{longtable}{|m{2.2cm}|m{2.4cm}|m{2cm}|m{2.7cm}|m{2cm}|m{3cm}|}
\caption{Tabla comparativa de los proyectos con semejanza al prototipo.}
\label{table:proyec}
\\ \hline
\multirow{2}{*}{Proyecto} & \multicolumn{1}{c|}{Módulo 1} & \multicolumn{1}{c|}{Módulo 2} & \multicolumn{1}{c|}{Módulo 3} & \multicolumn{1}{c|}{Módulo 4} & \multirow{2}{*}{Tipo de relación} \\ \cline{2-5}
 & Sistema de Sensores & Análisis de datos & Base de datos & Página Web & \\ \hline
\ref{sec:In1} & \begin{tabular}[c]{@{}c@{}}\checkmark\\ 1)Sensor de\\ pH\\ 2)Sensor de \\ temperatura\\ 3)Arduino\end{tabular} & & & \begin{tabular}[c]{@{}c@{}}\checkmark\\ Node Red\end{tabular} & \begin{tabular}[c]{@{}c@{}}Tecnología\\ implementada\\ y\\ problemática\\ a resolver\end{tabular} \\ \hline
\ref{sec:In2} & & \begin{tabular}[c]{@{}c@{}}\checkmark\\ Modelo de\\ Lógica \\ difusa\end{tabular} & & & \begin{tabular}[c]{@{}c@{}}Técnicas de \\ análisis \\ implementada\\ y problemática a \\ resolver\end{tabular} \\ \hline
\ref{sec:In3} & \begin{tabular}[c]{@{}c@{}}\checkmark\\ 1)Sensor de \\ temperatura\\ 2)ATMega32\end{tabular} & \begin{tabular}[c]{@{}c@{}}\checkmark\\ Sistema \\ con lógica\\ difusa\end{tabular} & & & \begin{tabular}[c]{@{}c@{}}Tecnología y \\ técnicas de \\ análisis \\ implementada\end{tabular} \\ \hline
\ref{sec:Na1} & \begin{tabular}[c]{@{}c@{}}\checkmark\\ 1)Sensor de\\ pH\\ 2)Raspberry Pi\\ Zero\end{tabular} & \begin{tabular}[c]{@{}c@{}}\checkmark\\ Red\\ Neuronal\\ NARX\end{tabular} & \begin{tabular}[c]{@{}c@{}}\checkmark\\ 1)Influxdb\\ cloud\\ 2)SQL\end{tabular} & \begin{tabular}[c]{@{}c@{}}\checkmark\\ Grafana\end{tabular} & \begin{tabular}[c]{@{}c@{}}Tecnología y\\ técnicas de \\ análisis\\ implementada\end{tabular} \\ \hline
\ref{sec:Na2} & \begin{tabular}[c]{@{}c@{}}\checkmark\\ 1)Sensor de\\ pH\\ 2)Sensor de \\ temperatura\\ 3)Arduino\\ Mega\end{tabular} & & & & \begin{tabular}[c]{@{}c@{}}Tecnología \\ implementada\end{tabular} \\ \hline
\ref{sec:Na3} & \begin{tabular}[c]{@{}c@{}}\checkmark\\ 1)Sensor de\\ $CO_{2}$\\ 3)Arduino\\ Nano\end{tabular} & & & \begin{tabular}[c]{@{}c@{}}\checkmark\\ LABVIEW\\ 2014\end{tabular} & \begin{tabular}[c]{@{}c@{}}Tecnología \\ implementada\end{tabular} \\ \hline
\ref{sec:Na4} & \begin{tabular}[c]{@{}c@{}}\checkmark\\ 1)Sensor de\\ pH\\ 2)Sensor de \\ temperatura\\ 3)Raspberry \\Pi 3\end{tabular} & & & & \begin{tabular}[c]{@{}c@{}}Tecnología \\ implementada\end{tabular} \\ \hline
\ref{sec:IPN1} & \begin{tabular}[c]{@{}c@{}}\checkmark\\ 1)Sensor de\\ pH\\ 2)Sensor de \\ temperatura\\ 3)ESP32\end{tabular} & & \begin{tabular}[c]{@{}c@{}}\checkmark\\ Servidor de \\Azure\\ Database\\ para \\MySQL\end{tabular} & \begin{tabular}[c]{@{}c@{}}\checkmark\\ Framework\\ express\end{tabular} & \begin{tabular}[c]{@{}c@{}}Tecnología \\ implementada\end{tabular} \\ \hline
\ref{sec:IPN2} & \begin{tabular}[c]{@{}c@{}}\checkmark\\ 1)Sensor de\\ $CO_{2}$\end{tabular} & \checkmark & \begin{tabular}[c]{@{}c@{}}\checkmark\\ AWS \\ Dynamo DB\end{tabular} & \checkmark & \begin{tabular}[c]{@{}c@{}}Tecnología \\ implementada\end{tabular} \\ \hline
\end{longtable}
\endgroup
\vspace{0.5cm}

Como se observa en la tabla anterior, los trabajos revisados implementan diversos
enfoques de monitoreo y análisis de datos aplicados a la producción de alimentos y
bebidas fermentadas, lo que evidencia tanto la complejidad como la relevancia de
estos sistemas. Cada propuesta fue diseñada con características particulares
adaptadas al medio y producto específico a monitorear. Sin embargo, aunque varios
de estos trabajos comparten una problemática similar a la abordada en el presente
proyecto, como los proyectos  \ref{sec:In1} y  \ref{sec:In2}, se identifican limitaciones importantes
en el desarrollo integral de cada módulo propuesto. Por otro lado, aunque el
proyecto  \ref{sec:Na1} contempla módulos similares, su enfoque está dirigido a una
problemática distinta. Esto resalta la necesidad de desarrollar un sistema completo,
con todos los módulos interconectados, orientado específicamente al monitoreo de
la fermentación artesanal del kéfir de leche, con el fin de mejorar la precisión,
trazabilidad y control del proceso